\section{两段燃烧规划与对比算法}
\subsection{控制参数化与优化目标}
采用两段燃烧及中间滑行结构，将连续控制函数压缩为十维设计向量：
\begin{equation}z=(\delta_1,\delta_2,\delta_3,\delta_4,\delta_5,q_1,q_2,t_1,t_c,t_2)^{\mathsf T}.\end{equation}
其中，前五项描述偏角节点，其余变量给出两段平台流率及燃烧、滑行时长。第一段偏角节点位于$0.2t_1$、$0.6t_1$和关机过渡起点，第二段偏角由$\delta_4$逐步过渡至$\delta_5$。节点之间采用线性插值，滑行段流率严格为零，启停过程保留有限时长过渡。

在$\xi\in[0,1]^{10}$搜索，$z=z_{\rm lb}+\xi\odot(z_{\rm ub}-z_{\rm lb})$。边界为
\begin{equation}
\begin{aligned}
z_{\rm lb}&=(0,-15,-15,-15,-15,180,120,80,20,25),\\
z_{\rm ub}&=(15,15,15,15,15,350,350,240,400,150).
\end{aligned}
\end{equation}
角度、流率和时间单位分别为度、kg/s和s。规划要求终端状态接近目标圆轨道，不限定注入位置的极角，采用以下归一化误差目标：
\begin{equation}
J(z)=\left(\frac{h_f-h_*}{10^4\,{\rm m}}\right)^2+
\left(\frac{u_f}{10^3\,{\rm m/s}}\right)^2+
\left(\frac{w_f-v_c}{10^3\,{\rm m/s}}\right)^2.\label{eq:J}
\end{equation}
燃料超限返回$10^3+[(M_{\rm used}-M_{\rm budget})/(10^3\,{\rm kg})]^2$；触地则加$100+(t_f-t_{\rm impact})/(1\,{\rm s})$。目标不含燃料、时间、动压或过载最小项，因此不同解的资源差异不代表算法直接优化了这些指标。

\subsection{粒子群优化}
PSO按个体最佳$p_i$和群体最佳$g$更新\cite{kennedy1995pso}：
\begin{equation}
\begin{aligned}
\nu_i^{k+1}&=\omega_k\nu_i^k+c_1r_1\odot(p_i-\xi_i^k)+c_2r_2\odot(g-\xi_i^k),\\
\xi_i^{k+1}&=\mathcal R(\xi_i^k+\nu_i^{k+1}).
\end{aligned}
\end{equation}
$\mathcal R$反射修复边界。采用全局最佳拓扑，$\omega_k$由1.1线性降至0.1，$c_1=c_2=1.49$，速度限制$[-0.2,0.2]$。上述设置在统一评价预算下执行，以比较群体搜索对终端误差的降低能力。

\subsection{实数编码遗传算法}
GA采用实数染色体、三元锦标赛、混合算术交叉和高斯变异\cite{holland1992}。父代$a,b$生成
\begin{equation}\xi'=\gamma\odot a+(1-\gamma)\odot b,\quad\gamma_\ell\sim U(-0.25,1.25).\end{equation}
交叉后，各坐标以$1/10$的概率增加标准差为0.1的高斯扰动，并通过反射处理越界。每代保留一个历史最优个体，以维持已获得的搜索精度；该精英的重新评价计入总预算。

\subsection{多起点模拟退火}
SA设置96条独立链，每轮各链扰动一个随机坐标，按Metropolis规则接受\cite{kirkpatrick1983}：
\begin{equation}P_{\rm accept}=\min\{1,\exp[-(J(\xi')-J(\xi))/T_k]\}.\end{equation}
温度由1按几何序列降至$10^{-6}$，提议标准差由0.2线性降至0.005。各链独立更新、不交换状态，最终返回全部访问点中的历史最优解。该设置将相同初始候选用于多起点搜索，与单链退火及并行回火分别对应不同的实验协议。

\subsection{等预算配对设计}
每方法使用96个候选（SA为96条链）、180轮更新，总评价$96(180+1)=17376$次。种子161--165形成五组配对实验；组内初始点逐元素一致，一半均匀随机，一半按两段燃烧燃料分配启发式生成。后续算法使用独立可复现随机流，规划扰动固定为20261020。

各次运行均在固定预算耗尽后终止，并将搜索历史最优点提交至两档高精度复核。复核仅评价固定控制，不引入局部最小化或参数再搜索。性能指标包括完整收敛曲线、终端残差、目标中位数、四分位范围及验收通过数。配对种子用于观察优化器随机性，而实际扰动环境保持一致。
